\begin{abstract}
\noindent
Threshold rules promise that a resource remains above a limit, but simulation-based assessments usually estimate survival probabilities rather than establish a guarantee under a specified model class. We derive closed-loop robust viability kernels for one-dimensional monotone maps and apply the results, under protocols frozen on 2026-08-26, to Northern cod (NAFO 2J3KL; limit reference point $\theta=884.6$ kt) and Edwards Aquifer J-17 (triggers at 618 and 660 ft). For a continuous, strictly increasing map with a non-binding ceiling, the finite-horizon kernel has a computable lower boundary; over increasing horizons, the kernels are either the whole safe set, converge to a half-line above the smallest fixed point, or become empty at a finite horizon. For the Schaefer cod map, the regimes are governed by $C^*=g(\theta)-|e|$ and $C_{\mathrm{vac}}=g_{\max}-|e|$; for the affine aquifer map, by its worst-floor attractor. A staged rule affects protection at a threshold only through the control active at that threshold.

For cod, the q10 floor supports a constructive bound of $91.6$ kt, whereas the q05 and worst floors do not make the LRP self-viable. Exact bounded-defect horizons range from 8--13 years under the worst and q05 floors and 12--28 years under q10; the registered Lipschitz certificate gives 5--7 years. For the aquifer, a flat pumping cut of at least 7.2\% secures 618 ft; on observed heads the Stage-I rule supplies 3.3\% more than a flat 90\% cap at matched protection. No admissible rule secures 660 ft, even at zero pumping. At 618 ft, the declared-defect horizon is three years for positive-pumping rules and four at zero pumping; the out-of-sample defect reduces these to two and three years, respectively. The exact recursion extends cod horizons but leaves aquifer horizons unchanged, while recording boundary departures for staged controls and for zero pumping at the shortest horizons. Neither system admits an indefinite certificate under its declared defect. The framework links the catch-sensitive cod horizon to curvature away from the LRP and the aquifer horizon to the defect-to-contraction margin against its ceiling. An independent core reproduces 926 comparisons with the registered campaign artefacts; the data, protocols, code and verification record are distributed with the manuscript.
\end{abstract}

\noindent\textbf{Keywords:} viability kernel; robust invariance; limit reference point; critical period management; Northern cod; Edwards Aquifer; certification; model defect; reassessment cadence.

\section{Introduction}\label{sec:intro}

A threshold rule promises that a state variable will be held above a line. The 2016 Northern cod assessment fixed the limit reference point (LRP) at the 1983--1989 mean spawning stock biomass, $\LRP=884.6$ kt, and the Canadian precautionary framework attaches catch consequences to crossing it \cite{DFO2009,DFO2016}. The Edwards Aquifer Authority manages pumping through staged cuts triggered when the J-17 index well falls below 660, 650, 640 and 630 ft above mean sea level, with 618 ft the elevation below which spring flows and the endangered-species commitments of the Habitat Conservation Plan are in jeopardy \cite{EAA2024,EARIP2021,NRC2015}. In both cases the rule is stated as a guarantee, but the evidence offered for it is probabilistic: stochastic projections, management strategy evaluation, scenario ensembles \cite{Butterworth2007,Punt2016}. Those tools estimate how often the threshold holds. They do not answer the question the rule poses, which is whether the threshold can be held against everything a fitted model admits---whether it can be \emph{certified}.

Viability theory supplies the right object for that question. The viability kernel of a constraint set under a controlled dynamical system is the set of states from which the constraint can be maintained for ever \cite{Aubin1991}; its finite-horizon and robust variants are the states from which it can be maintained for $T$ steps, or against a class of disturbances \cite{BertsekasRhodes1971,RegnierDeLara2015}. The theory has been applied to fisheries and groundwater \cite{Bene2001,Martinet2007,Doyen2012,Krawczyk2013,Pereau2018,Pereau2019,Oubraham2018}, usually with numerically approximated kernels in two or more dimensions. The present paper takes a different route: it stays in one dimension, where the kernel of a threshold under a monotone map can be written down, and uses the closed form to say exactly what a frozen model does and does not establish about a rule.

Two systems are treated under one frame. They were chosen for contrast rather than similarity. The cod map is expansive at its reference point---the slope of the one-step map at the LRP is $F'(\LRP)=1.153$---so disturbances are amplified near the line the rule is meant to protect; the aquifer map is a contraction ($a=0.746$), so disturbances decay. The cod residuals are dominated by the 1991--1992 collapse years, which are seven to ten residual standard deviations from zero; the aquifer residuals are small in sample and larger out of sample. The cod safe set has no upper limit that matters; the aquifer has a ceiling at 710 ft that the backward recursion runs into. Everything that follows was computed under protocols frozen before the computations were made (both on 2026-08-26); the protocols are reproduced in the supplement.

The paper makes four contributions.

\begin{enumerate}[leftmargin=2em]
\item \emph{A closed-form certification frame for one-dimensional threshold management} (Section~\ref{sec:frame}). Theorem~\ref{thm:kernel} characterises the $T$-step and infinite-horizon closed-loop kernels of a threshold under any strictly increasing map with a non-binding ceiling: the kernel is the whole safe set, a half-line above the smallest fixed point, or empty beyond an explicit finite horizon. For a surplus-production stock the characterisation collapses to two constants, $\Cstar=g(\LRP)-|e|$ and $\Cvac=\gmax-|e|$ (Corollary~\ref{cor:cod}); for an affine aquifer map to one attractor and a closed-form emptying horizon (Corollary~\ref{cor:edw}). Proposition~\ref{prop:stage} shows that a staged rule acts on a threshold only through the stage active at that threshold.
\item \emph{An exact treatment of model defect} (Theorem~\ref{thm:exact}, Propositions~\ref{prop:lip}--\ref{prop:finite}). The registered protocols certified against a uniform model defect $\varepsilon$ by inflating the threshold with a Lipschitz accumulation $r_T$. We show that the exact robust kernel of the bounded-defect class is computable by the same interval recursion, that the Lipschitz certificate coincides with it for affine maps wherever the recursion binds, and that it is conservative for maps whose slope falls away from the threshold. For cod the exact kernel extends the horizon and makes it catch-sensitive; for the aquifer it leaves the horizon unchanged, although boundaries depart from the Lipschitz certificate at zero pumping for $T=1,2$ and under the staged rules. In both systems the certified horizon is finite for every admissible rule, and the condition for an indefinite certificate is stated and shown to fail.
\item \emph{Resolved answers for two managed systems} (Sections~\ref{sec:cod}--\ref{sec:edwards}, Table~\ref{tab:resolved}). For cod: the LRP is self-viable under constant catches up to $91.6$ kt under the 10th-percentile residual floor and under no catch at all under the 5th-percentile or worst floors; the Stage-I rule, the declared cascade and a flat 60-kt cap have identical kernels; the constructive bound is identified from the data even though the carrying capacity is not. For the aquifer: 618 ft is secured by any flat cut of at least $7.2\%$; the Stage-I rule is retained against the flat family because it ties the 90\% cap in protection and supplies $3.3\%$ more water; the 660-ft trigger cannot be held from itself under any recharge floor even with zero pumping; and the 2004 out-of-sample residual exceeds the declared defect, reducing the horizon for positive-pumping rules from three to two years while zero pumping retains three.
\item \emph{A cross-system reading of the certified horizon} (Section~\ref{sec:synthesis}). The horizons are finite but differ in their margins: under the cod Lipschitz certificate they are six years under the worst/q05 floors and seven under q10 for catches up to 150 kt, while the exact kernels extend to 8--13 and 12--28 years, respectively. For the aquifer the declared-defect horizon is three years for positive-pumping rules (four at zero pumping), and the out-of-sample horizon is two/three. The six-year International Whaling Commission Implementation Review cycle matches the registered cod certificate for catches up to 180 kt under the worst/q05 floors; the aquifer result under the out-of-sample defect calls for a shorter two-year reassessment interval.
\end{enumerate}

The paper is organised as follows. Section~\ref{sec:frame} develops the frame. Sections~\ref{sec:cod} and \ref{sec:edwards} apply it to cod and to the aquifer, each with its stochastic and identification layers. Section~\ref{sec:synthesis} reads the two systems together. Section~\ref{sec:discussion} states limitations, Section~\ref{sec:verification} the verification record, Section~\ref{sec:related} the relation to prior work, and Section~\ref{sec:conclusion} concludes. The supplement (S1--S14) carries the data, the full kernel tables, the stochastic and bootstrap grids, the Edwards replays and the symbol list.
